%% ma: L10-B15-0110
%% lop: 10
%% chuong: 6
%% bai: 15
%% dang: 10.15.5
%% loai: TLN
%% muc: TH
%% dap_an: "20"
%% nguon: "Ảnh giáo viên gửi 25/11/2026 (lô 2): Bài tập trắc nghiệm – Toán 10 Bài 15 Hàm số (sự biến thiên, đồ thị), Phần III Câu 2"
%% kien_thuc: "Lớp 10 Bài 15 (đồng biến, nghịch biến); hàm số bậc nhất y=ax+b đồng biến khi a>0 (lớp 9)"
%% trang_thai: chinh-thuc
%% ghi_chu: "Biến thể thay số của L10-B15-0048."
\begin{ex}
Có bao nhiêu giá trị nguyên của tham số $m$ trong đoạn $[-20;20]$ để hàm số $f(x)=(2m-1)x+m+3$ đồng biến trên $\mathbb R$?
\shortans{20}
\loigiai{Cần $2m-1>0\Leftrightarrow m>\tfrac12$. Các số nguyên trong $[-20;20]$: $m=1;2;\dots;20$. Có $20$ giá trị.}
\end{ex}
